\section{Discussion}
\label{sec:discussion}

\subsection{A coordination failure we caused ourselves}
\label{sec:discussion-deadlock}

The most instructive event in this work was not planned. Our launcher---the
subagent API we used to start ranks---turned out to have a concurrency limit
of ten. We asked for twelve translator ranks; nine started, three were
refused, and the nine that started blocked in the exclusive scan waiting for
peers that did not exist. Because every one of the nine held a launcher
slot while blocked, no slot ever freed, and the three missing ranks could
never start. The job was deadlocked by a resource limit two layers below
the protocol, and no rank could observe the cause: each was correctly
waiting for a message that a peer would have sent had it been allowed to
run.

This is a textbook resource deadlock, and it is worth reporting for three
reasons. First, it is exactly the class of failure a harness author does not
anticipate, because the launcher's capacity is not part of the program.
Second, it is invisible without a global view: \code{ampi peers} showed
three ranks that had never checked in, which is the whole diagnosis, and
nothing in any rank's local state would have revealed it. Third, the
protocol already contained the correct response---\op{Comm\_shrink} over the
ranks that did arrive---and nothing that would notice the situation in
useful time.

That gap is now closed, and closing it required separating two conditions
that had been conflated. A rank that registered and then went quiet may come
back, so it gets the failure timeout, which must be generous because an
agent's turn is long. A rank that has \emph{never registered} was never
launched; no amount of waiting will produce it, and spending the failure
timeout on it is pure loss. The runtime now tracks registration as a durable
fact separate from liveness and applies a short \emph{roll-call} deadline,
measured from the run's creation rather than from the observing rank's start
so that a late joiner does not restart everyone's clock. In the incident
above, eight ranks blocked for twenty minutes on a broadcast whose root had
not started; with the roll call they would have been told within seconds
which ranks were missing and that a shrink was required. We report this as a
fix prompted by the failure rather than a design we anticipated.

The general lesson is that a collective's all-participants requirement makes
the launcher's capacity a correctness property, not a performance one. MPI
gets this for free because \code{mpiexec} either allocates $p$ slots or
fails. An agent launcher that partially fulfils a request produces a job
that starts, looks healthy, and never finishes.

\subsection{Agents disobey the protocol, and the protocol must survive it}
\label{sec:discussion-disobedience}

We observed two protocol violations by agent ranks in a single run, and
neither was a model quality problem in the usual sense.

Rank 1 decided that receiving its work assignment through \code{ampi
scatter} was unnecessary. It read the corpus from disk instead and inferred
which chunk was its own. It inferred wrongly---the scatter had assigned it
chapter~1, and it took chapter~2---and it translated the wrong chapter
competently and at length. Worse, skipping the collective left its
collective counter one behind every peer's, so from then on it would have
labelled every message with a tag nobody was listening for. This is the
failure that motivated collective-order verification
(\Cref{sec:design}); with that mechanism the run would have reported
\emph{``collective \#2 was issued as `scatterv' here and as `exscan' by
ranks [2,4]; this rank skipped a collective''} within seconds instead of
hanging.

Rank 3 reordered its program, translating before running the scan it was
supposed to run first, and then stopped without completing its remaining
collectives. Its translation was good. Its participation was not.

Both are instances of what the MAST taxonomy calls disobeying the task
specification~\cite{cemri2025mast}, and both were locally reasonable: an
agent that can see the corpus has no obvious reason to wait for a message
containing part of it. The design conclusion we draw is that a protocol for
agents cannot rely on participants following it, and must therefore make
deviation \emph{detectable} rather than merely forbidden. Contract checking
does this for message content; collective-order verification does it for
control flow; the failure lattice does it for liveness. We do not claim the
set is complete.

There is a second, more uncomfortable conclusion, and we have direct
evidence for it. An agent that can reach shared state by a path the protocol
does not mediate will sometimes take it, and the protocol's guarantees
evaporate silently when it does. MPI does not have this problem, because a
rank physically cannot read another rank's address space.

Our own software harness supplied the counter-example. Its rank program told
every rank to publish its interface with \code{--file iface.json} and did
not give the ranks separate working directories, so eight agents shared one
directory and one filename. Each rank wrote \code{iface.json} and then asked
the runtime to publish it; whichever rank published second published its
neighbour's file. The symptom, reported by the rank that noticed, was that
the interface board appeared to have \emph{aliased} two keys, with
byte-identical previews and identical token counts --- which is exactly what
the window faithfully storing the same bytes twice looks like from the
outside. We pinned window isolation with a test before believing the report,
and the window is correct: the corruption happened entirely outside the
protocol, in the one plane the protocol does not mediate.

The lesson is sharper than ``sandbox your agents''. Every AgentMPI operation
that names a file is a place where the protocol's isolation stops and the
filesystem's does not begin. Our translation program gave each rank a
private scratch directory and had no such failure; the software program did
not and did. The remedy we would now insist on is that a launcher assign
each rank a private working directory and that rank programs never name a
bare relative path --- and, in a serious deployment, that ranks run in
sandboxes whose only shared view is the run directory.

\subsection{Four bugs, and what they have in common}
\label{sec:discussion-bugs}

The experiments found four defects in our own implementation, and they
share a shape worth naming.

\emph{Sub-communicator misrouting.} Envelopes carried communicator-local
ranks; the device addresses inboxes by world rank. On \op{COMM\_WORLD} the
two numberings coincide, so 260 tests passed and every collective on a
shrunk or split communicator was misdelivered.

\emph{Non-durable collective counters.} The counter separating one
collective's traffic from the next was persisted at process exit. A rank
killed mid-collective---an agent's shell timing out---reused the number, and
was permanently one behind its peers.

\emph{Consumption at poll time.} Our first matching engine marked a message
consumed when it was polled rather than when it was matched, so a rank
implemented as a sequence of short-lived commands destroyed every message it
polled but did not match.

\emph{No run identity.} We deleted a wedged run and created a new one at the
same path while ranks of the old one were still blocked. Those ranks
re-attached, inherited the previous run's collective counter from their own
durable state, numbered every collective one too high, and matched the wrong
message: a rank's broadcast returned the payload of the coordinator's
scatter, and its scatter returned the integer a barrier had sent, which
surfaced as a type error a long way from its cause. An epoch counter guards
against stale messages \emph{within} a run and is silent about a recycled
directory. Durable state now carries the run identity from the manifest and
is discarded when it does not match, and a rank whose run directory has
vanished reports \code{ERR\_NO\_RUN} and stops rather than issuing further
collectives.

All four are the same mistake: an assumption that holds when a rank is one
long-lived process on the world communicator inside a job it cannot outlive,
and fails when it is a succession of short-lived processes, on a communicator
that can change, in a run whose transport can be deleted underneath it. MPI
never has to consider the fourth, because a job's transport is process-scoped
and dies with the job; a directory does not. The generalisable lesson is that
\emph{the endpoint being stateless is a load-bearing property of an agent
protocol}, not an implementation convenience, and every piece of protocol
state must be examined against it --- including the question of which run the
state belongs to. A test suite that only exercises \op{COMM\_WORLD} with
in-process ranks in a run that outlives them will not find any of the four.

It is worth recording how the fourth was diagnosed, because it bears on
\Cref{sec:discussion-disobedience}. The affected rank worked out the cause
itself --- that its persisted counter was one ahead, that this was why its
broadcast had returned the wrong payload --- observed that a second process
was concurrently acting as the same rank, and then \emph{declined to issue
its remaining collectives} on the grounds that doing so would overwrite state
a live peer depended on. That is the behaviour one wants and cannot rely on.
The protocol should have made the situation impossible rather than merely
diagnosable, which is what the run-identity check now does.

The second bug also exposes a limit of our own diagnostics. Collective-order
verification compares operation \emph{names} at a sequence number, so it
detects a rank that skipped a collective and not a rank that repeated one:
a replayed barrier has the same name as the barrier its peers ran. Detecting
repetition would need a per-invocation nonce that all ranks derive
identically, which we have not designed.

\subsection{Operator algebra deserves more attention than it gets}

The result we would most like other designers to take away is the one about
commutativity (\Cref{sec:ops}). Every framework that merges contributions
from several agents is applying a reduction operator, and almost none of
them ask whether the operator is associative or commutative. The
consequences are not subtle: our own first implementation of a
shared-terminology scan used a union, which is commutative, and it left the
participants disagreeing about a decision the scan existed to make. The
error was invisible in small runs and appeared as an 18-point quality drop
at twelve ranks.

Commutativity is the wrong default for consensus, and precedence is
inexpressible without a rank order. We suspect this generalises well beyond
terminology to any convention an agent team must settle---schema shapes,
naming, file layout, API contracts---and that a good deal of reported
multi-agent inconsistency is this bug.

\subsection{Limitations}

\paragraph{Scale.} Our agent runs are at 12 ranks and our software build at
8. The microbenchmarks reach 128 ranks, but with scripted executors: they
measure the protocol, not the agents. We cannot claim from this data that
agent \emph{quality} holds at 128 ranks, and we would expect it not to for
allgather-shaped workloads for the reasons in \Cref{sec:cost}.

\paragraph{One model family.} All agent ranks were the same class of coding
agent. \ampi{}'s heterogeneity machinery---\code{RankSpec},
\op{Comm\_split\_type} by model, cost-aware selection---is implemented and
tested but not exercised against genuinely different models.

\paragraph{Statistical power.} We report single runs of each configuration.
Agent runs are expensive and slow, and turn durations are heavy-tailed, so
the wall-clock numbers in particular should be read as one draw rather than
as an estimate of a mean. The round counts are exact and deterministic; the
consistency metrics are deterministic given the outputs; only the timings
are noisy.

\paragraph{Operator intervention.} Two ranks in the glossary run required
manual completion of protocol steps they had skipped
(\Cref{sec:discussion-disobedience}). The translations they contributed are
their own work, but the run is not a clean unattended execution, and we
report it as such.

\paragraph{The failure detector is weak for shell-mode ranks.} A rank
implemented as a sequence of short-lived commands cannot heartbeat while it
thinks, so its liveness timeout must exceed its longest turn. In our runs
that meant a 30-minute timeout, which is a long time to wait to discover
that an agent has died. A supervisor process per rank would fix this and is
the obvious next step.

\paragraph{Not evaluated against alternative frameworks.} We compare against
an unstructured baseline, not against LangGraph or AutoGen. A fair
comparison would require reimplementing the same workload in each, and the
result would largely measure how much of the protocol those frameworks
oblige the author to rebuild by hand---which is our claim, but not one this
paper's data establishes.
